% !TEX program = xelatex
% ============================================================
%  نشرة إخبارية (Newsletter) Template
%  قالب النشرة الإخبارية مع ترويسة ملوّنة وقصص رئيسية وثانوية
%  Compile with: xelatex main.tex (run twice)
% ============================================================
%  Features:
%    - Colored masthead bar (newsletter name + issue + date)
%    - Lead story (headline + 2 paragraphs + figure placeholder)
%    - Two secondary stories in two-column layout
%    - Upcoming events box (tcolorbox list)
%    - Quick links section
%    - Footer (subscribe + contact)
% ============================================================

\documentclass[11pt,a4paper]{article}

\usepackage[a4paper,margin=2cm,top=2cm,bottom=2cm,headheight=16pt]{geometry}
\usepackage{graphicx}
\usepackage{array}
\usepackage{booktabs}
\usepackage{tabularx}
\usepackage{multirow}
\usepackage{enumitem}
\usepackage{float}
\usepackage{titlesec}
\usepackage{fancyhdr}
\usepackage{setspace}
\usepackage{xcolor}
\usepackage{tikz}
\usepackage{tcolorbox}
\tcbuselibrary{skins}
\usepackage{multicol}

% ---- Arabic/RTL packages (this order, before hyperref) ----
\usepackage{bidi}
\usepackage[no-math]{fontspec}
\usepackage[quiet]{polyglossia}

\setmainfont[Scale=1]{NotoKufiArabic-Regular.ttf}
\newfontfamily\arabicfont[Script=Arabic,Scale=0.95]{NotoKufiArabic-Regular.ttf}
\newfontfamily\englishfont[Scale=0.95]{TeX Gyre Termes}

\setdefaultlanguage[calendar=gregorian,hijricorrection=1]{arabic}
\setotherlanguage[variant=british]{english}

\usepackage[breaklinks,hidelinks]{hyperref}
\hypersetup{pdftitle={نشرة إخبارية},pdfauthor={اسم الشركة}}

% ============================================================
% Accent color
% ============================================================
\definecolor{accent}{HTML}{2E5A88}
\definecolor{accentdark}{HTML}{1F3A5F}
\definecolor{accentlight}{HTML}{E8EFF5}
\definecolor{eventbg}{HTML}{FFF8E1}
\definecolor{eventrule}{HTML}{D4A017}
\definecolor{storybg}{HTML}{F8F8F8}

% ============================================================
% Section formatting
% ============================================================
\titleformat{\section}{\large\bfseries\color{accentdark}}{\thesection}{0.5em}{}
\titlespacing*{\section}{0pt}{0.8ex plus 0.2ex}{0.5ex plus 0.1ex}

% ============================================================
% Header / footer
% ============================================================
\pagestyle{fancy}
\fancyhf{}
\fancyhead[R]{\small\itshape نشرة إخبارية (\LR{Newsletter})}
\fancyhead[L]{\small اسم الشركة}
\fancyfoot[C]{\small\thepage}
\renewcommand{\headrulewidth}{0pt}

% ============================================================
% Custom tcolorbox environments
% ============================================================
% Upcoming events box
\newtcolorbox{eventsbox}{
  enhanced,
  colback=eventbg,
  colframe=eventrule,
  boxrule=0.8pt,
  arc=3pt,
  left=10pt,right=10pt,top=8pt,bottom=8pt,
  title={\bfseries الفعاليات القادمة},
  colbacktitle=eventrule,
  coltitle=white,
  fonttitle=\bfseries,
}

% Story box for secondary stories
\newtcolorbox{storybox}{
  enhanced,
  colback=storybg,
  colframe=accentlight,
  boxrule=0.6pt,
  arc=2pt,
  left=8pt,right=8pt,top=6pt,bottom=6pt,
}

\setstretch{1.3}
\setlength{\parindent}{0pt}
\setlength{\parskip}{0.3em}

\begin{document}
\thispagestyle{fancy}

% ============================================================
% MASTHEAD (colored bar)
% ============================================================
\noindent\colorbox{accentdark}{\parbox{\dimexpr\textwidth-2\fboxsep\relax}{%
\centering
\vspace{0.3em}
{\color{white}\Huge\bfseries النشرة الإخبارية}\\[0.2em]
{\color{white}\large العدد --- \quad|\quad \today}\\[0.2em]
{\color{white}\small\itshape اسم الشركة \textemdash{} نشرة شهرية}
\vspace{0.2em}
}}

\vspace{0.6em}

% ============================================================
% LEAD STORY
% ============================================================
% TODO: استبدل عنوان ونص القصة الرئيسية.
{\LARGE\bfseries\color{accentdark} عنوان القصة الرئيسية هنا}

\vspace{0.3em}
{\small\itshape \LR{By} اسم الكاتب \quad|\quad \today}

\vspace{0.4em}
% Figure placeholder
\begin{center}
\fbox{\parbox[c][4cm][c]{0.85\textwidth}{\centering\large مكان الصورة\\(\LR{Figure placeholder})}}
\end{center}

\vspace{0.4em}
% TODO: اكتب فقرتي القصة الرئيسية.
هذه هي الفقرة الأولى من القصة الرئيسية. اكتب هنا ملخصاً%
جذاباً للخبر الأهم في هذا العدد، مع التركيز على%
---من، ماذا، لماذا--- في الجمل الأولى.

هذه هي الفقرة الثانية من القصة الرئيسية. وسّع التفاصيل هنا،%
وأضف الاقتباسات أو البيانات الداعمة. يمكن استخدام%
\LR{\textbf{Read more \textgreater{}}} للربط بالنص الكامل على الموقع.

\vspace{0.3em}\hrule\vspace{0.6em}

% ============================================================
% SECONDARY STORIES (two-column)
% ============================================================
\begin{multicols}{2}
\setlength{\columnsep}{1.5em}
\raggedright

% ---- Secondary story 1 ----
\begin{storybox}
{\large\bfseries\color{accentdark} عنوان القصة الثانية}\\[0.2em]
{\small\itshape \today}\\[0.3em]
% TODO: اكتب نص القصة الثانية.
وصف مختصر للقصة الثانية. اكتب فقرة أو فقرتين%
تلخصان الخبر أو الموضوع المهم.
\end{storybox}

\vspace{0.4em}

% ---- Secondary story 2 ----
\begin{storybox}
{\large\bfseries\color{accentdark} عنوان القصة الثالثة}\\[0.2em]
{\small\itshape \today}\\[0.3em]
% TODO: اكتب نص القصة الثالثة.
وصف مختصر للقصة الثالثة. اكتب فقرة أو فقرتين%
تلخصان الخبر أو الموضوع المهم.
\end{storybox}

\end{multicols}

\vspace{0.5em}

% ============================================================
% UPCOMING EVENTS BOX
% ============================================================
% TODO: استبدل بقائمة الفعاليات الفعلية.
\begin{eventsbox}
\begin{itemize}[topsep=0.2em,itemsep=0.25em,leftmargin=1.4em]
    \item \textbf{\today} \textemdash{} ---اسم الفعالية---، ---المدينة---
    \item \textbf{---التاريخ---} \textemdash{} ---اسم الفعالية---، ---المدينة---
    \item \textbf{---التاريخ---} \textemdash{} ---اسم الفعالية---، ---المدينة---
\end{itemize}
\end{eventsbox}

\vspace{0.5em}

% ============================================================
% QUICK LINKS
% ============================================================
\section{روابط سريعة}
% TODO: استبدل بالروابط الفعلية.
\begin{itemize}[topsep=0.2em,itemsep=0.15em,leftmargin=1.4em]
    \item \LR{\href{https://www.example.com}{الموقع الإلكتروني}}
    \item \LR{\href{https://www.example.com/blog}{المدونة}}
    \item \LR{\href{https://www.example.com/products}{المنتجات}}
    \item \LR{\href{https://www.example.com/support}{الدعم الفني}}
\end{itemize}

\vspace{0.5em}

% ============================================================
% FOOTER (subscribe + contact)
% ============================================================
\noindent\colorbox{accentlight}{\parbox{\dimexpr\textwidth-2\fboxsep\relax}{%
\vspace{0.4em}
\begin{center}
\textbf{اشترك في نشرتنا الإخبارية} \textemdash{}%
أرسل بريداً إلى%
\LR{\href{mailto:newsletter@example.com}{newsletter@example.com}}%
\quad|\quad%
تواصل معنا: \LR{info@example.com}
\end{center}
\vspace{0.3em}
{\centering\small\itshape \LR{www.example.com} \quad\textbullet\quad اسم الشركة \quad\textbullet\quad جميع الحقوق محفوظة \today\par}
\vspace{0.3em}
}}

\end{document}
